\documentclass{amsart}
% For dealing with references we use the comment environment
\usepackage{verbatim}
\newenvironment{reference}{\comment}{\endcomment}
%\newenvironment{reference}{}{}
\newenvironment{slogan}{\comment}{\endcomment}
\newenvironment{history}{\comment}{\endcomment}

% For commutative diagrams we use Xy-pic
\usepackage[all]{xy}

% We use 2cell for 2-commutative diagrams.
\xyoption{2cell}
\UseAllTwocells

% We use multicol for the list of chapters between chapters
\usepackage{multicol}

% This is generally recommended for better output
\usepackage{lmodern}
\usepackage[T1]{fontenc}

% For cross-file-references

% Package for hypertext links:
\usepackage{hyperref}

% For any local file, say "hello.tex" you want to link to please

% Theorem environments.
%
\theoremstyle{plain}
\newtheorem{theorem}[subsection]{Theorem}
\newtheorem{proposition}[subsection]{Proposition}
\newtheorem{lemma}[subsection]{Lemma}

\theoremstyle{definition}
\newtheorem{definition}[subsection]{Definition}
\newtheorem{example}[subsection]{Example}
\newtheorem{exercise}[subsection]{Exercise}
\newtheorem{situation}[subsection]{Situation}

\theoremstyle{remark}
\newtheorem{remark}[subsection]{Remark}
\newtheorem{remarks}[subsection]{Remarks}

\numberwithin{equation}{subsection}

% Macros
%
\def\lim{\mathop{\mathrm{lim}}\nolimits}
\def\colim{\mathop{\mathrm{colim}}\nolimits}
\def\Spec{\mathop{\mathrm{Spec}}}
\def\Hom{\mathop{\mathrm{Hom}}\nolimits}
\def\Ext{\mathop{\mathrm{Ext}}\nolimits}
\def\SheafHom{\mathop{\mathcal{H}\!\mathit{om}}\nolimits}
\def\SheafExt{\mathop{\mathcal{E}\!\mathit{xt}}\nolimits}
\def\Sch{\mathit{Sch}}
\def\Mor{\mathop{\mathrm{Mor}}\nolimits}
\def\Ob{\mathop{\mathrm{Ob}}\nolimits}
\def\Sh{\mathop{\mathit{Sh}}\nolimits}
\def\NL{\mathop{N\!L}\nolimits}
\def\CH{\mathop{\mathrm{CH}}\nolimits}
\def\proetale{{pro\text{-}\acute{e}tale}}
\def\etale{{\acute{e}tale}}
\def\QCoh{\mathit{QCoh}}
\def\Ker{\mathop{\mathrm{Ker}}}
\def\Im{\mathop{\mathrm{Im}}}
\def\Coker{\mathop{\mathrm{Coker}}}
\def\Coim{\mathop{\mathrm{Coim}}}

% Boxtimes
%
\DeclareMathSymbol{\boxtimes}{\mathbin}{AMSa}{"02}

%
% Macros for moduli stacks/spaces
%
\def\QCohstack{\mathcal{QC}\!\mathit{oh}}
\def\Cohstack{\mathcal{C}\!\mathit{oh}}
\def\Spacesstack{\mathcal{S}\!\mathit{paces}}
\def\Quotfunctor{\mathrm{Quot}}
\def\Hilbfunctor{\mathrm{Hilb}}
\def\Curvesstack{\mathcal{C}\!\mathit{urves}}
\def\Polarizedstack{\mathcal{P}\!\mathit{olarized}}
\def\Complexesstack{\mathcal{C}\!\mathit{omplexes}}
% \Pic is the operator that assigns to X its picard group, usage \Pic(X)
% \Picardstack_{X/B} denotes the Picard stack of X over B
% \Picardfunctor_{X/B} denotes the Picard functor of X over B
\def\Pic{\mathop{\mathrm{Pic}}\nolimits}
\def\Picardstack{\mathcal{P}\!\mathit{ic}}
\def\Picardfunctor{\mathrm{Pic}}
\def\Deformationcategory{\mathcal{D}\!\mathit{ef}}

\setlength{\emergencystretch}{3em}
\begin{document}
\title[Stacks Project, Tag 0CSJ]{492 --- Derived Chow lemma for separated finitely presented schemes over arbitrary rings}
\maketitle
\noindent Copyright (C) 2005--2025 Johan de Jong. Stacks Project authors. Source: \href{https://stacks.math.columbia.edu/tag/0CSJ}{Tag 0CSJ}.

\noindent Editorial difficulty note (not author text). Difficulty H2: A graph closure over a Noetherian approximation constructs the geometric modification and a coherent supported object. Pulling this object back to the original base requires an explicit Tor-independence argument to retain its structure-sheaf identification on the prescribed projective open, while preserving pseudo-coherence and support..

\noindent Editorial provenance note (not author text). History: excerpt from revision \texttt{a0444\allowbreak 6e57e\allowbreak c1fbc\allowbreak 252a8\allowbreak 71afc\allowbreak ec775\allowbreak 2fb28\allowbreak 07b14}, more-morphisms.tex, lines 20336--20438. Reference presentation was adapted; original-author context and core support blocks were retained or added. The mathematical wording is unchanged. Licensed under GNU FDL 1.2 or later; no invariant sections or cover texts. See ../licenses/COPYING. Context blocks reproduce author definitions and supporting results; they are not additional counted problems.

\begin{lemma}[Derived Chow's lemma]
\label{lemma-derived-chow}
Let $A$ be a ring. Let $X$ be a separated scheme of finite presentation
over $A$. Let $x \in X$. Then there exist
an open neighbourhood $U \subset X$ of $x$,
an $n \geq 0$,
an open $V \subset \mathbf{P}^n_A$,
a closed subscheme $Z \subset X \times_A \mathbf{P}^n_A$,
a point $z \in Z$, and
an object $E$ in $D(\mathcal{O}_{X \times_A \mathbf{P}^n_A})$ such that
\begin{enumerate}
\item $Z \to X \times_A \mathbf{P}^n_A$ is of finite presentation,
\item $b : Z \to X$ is an isomorphism over $U$ and $b(z) = x$,
\item $c : Z \to \mathbf{P}^n_A$ is a closed immersion over $V$,
\item $b^{-1}(U) = c^{-1}(V)$, in particular $c(z) \in V$,
\item $E|_{X \times_A V} \cong
(b, c)_*\mathcal{O}_Z|_{X \times_A V}$,
\item $E$ is pseudo-coherent and supported on $Z$.
\end{enumerate}
\end{lemma}

\begin{proof}
We can find a finite type $\mathbf{Z}$-subalgebra $A' \subset A$
and a scheme $X'$ separated and of finite presentation over $A'$
whose base change to $A$ is $X$. See
Limits, Lemmas \href{https://stacks.math.columbia.edu/tag/01ZM}{lemma descend finite presentation (Tag 01ZM)} and
\href{https://stacks.math.columbia.edu/tag/01ZQ}{lemma descend separated finite presentation (Tag 01ZQ)}.
Let $x' \in X'$ be the image of $x$.
If we can prove the lemma for $x' \in X'/A'$, then
the lemma follows for $x \in X/A$.
Namely, if $U', n', V', Z', z', E'$ provide the solution
for $x' \in X'/A'$, then we can let
$U \subset X$ be the inverse image of $U'$,
let $n = n'$,
let $V \subset \mathbf{P}^n_A$ be the inverse image of $V'$,
let $Z \subset X \times \mathbf{P}^n$ be
the scheme theoretic inverse image of $Z'$,
let $z \in Z$ be the unique point mapping to $x$, and
let $E$ be the derived pullback of $E'$.
Observe that $E$ is pseudo-coherent by
Cohomology, Lemma \href{https://stacks.math.columbia.edu/tag/09U7}{lemma pseudo coherent pullback (Tag 09U7)}.
It only remains to check (5). To see this
set $W = b^{-1}(U) = c^{-1}(V)$ and $W' = (b')^{-1}(U) = (c')^{-1}(V')$ 
and consider the cartesian square
$$
\xymatrix{
W \ar[d]_{(b, c)} \ar[r] & W' \ar[d]^{(b', c')} \\
X \times_A V \ar[r] & X' \times_{A'} V'
}
$$
By Lemma \href{https://stacks.math.columbia.edu/tag/0CTA}{lemma case of tor independence (Tag 0CTA)} the schemes
$X \times_A V$ and $W'$ are Tor independent over $X' \times_{A'} V'$.
Hence the derived pullback of
$(b', c')_*\mathcal{O}_{W'}$ to $X \times_A V$
is $(b, c)_*\mathcal{O}_W$ by
Derived Categories of Schemes,
Lemma \href{https://stacks.math.columbia.edu/tag/08IB}{lemma compare base change (Tag 08IB)}.
This also uses that $R(b', c')_*\mathcal{O}_{Z'} = (b', c')_*\mathcal{O}_{Z'}$
because $(b', c')$ is a closed immersion and similarly for
$(b, c)_*\mathcal{O}_Z$.
Since $E'|_{U' \times_{A'} V'} =
(b', c')_*\mathcal{O}_{W'}$ we obtain
$E|_{U \times_A V} = (b, c)_*\mathcal{O}_W$
and (5) holds.
This reduces us to the situation described in the next
paragraph.

\medskip\noindent
Assume $A$ is of finite type over $\mathbf{Z}$.
Choose an affine open neighbourhood $U \subset X$ of $x$.
Then $U$ is of finite type over $A$.
Choose a closed immersion $U \to \mathbf{A}^n_A$ and denote
$j : U \to \mathbf{P}^n_A$ the immersion we get by composing
with the open immersion $\mathbf{A}^n_A \to \mathbf{P}^n_A$.
Let $Z$ be the scheme theoretic closure of
$$
(\text{id}_U, j) : U \longrightarrow X \times_A \mathbf{P}^n_A
$$
Since the projection $X \times \mathbf{P}^n \to X$ is separated,
we conclude from Morphisms, Lemma
\href{https://stacks.math.columbia.edu/tag/0CNG}{lemma scheme theoretic image of partial section (Tag 0CNG)}
that $b : Z \to X$ is an isomorphism over $U$.
Let $z \in Z$ be the unique point lying over $x$.

\medskip\noindent
Let $Y \subset \mathbf{P}^n_A$ be the scheme theoretic
closure of $j$. Then it is clear that $Z \subset X \times_A Y$
is the scheme theoretic closure of
$(\text{id}_U, j) : U \to X \times_A Y$.
As $X$ is separated, the morphism
$X \times_A Y \to Y$ is separated as well.
Hence we see that $Z \to Y$ is an isomorphism over
the open subscheme $j(U) \subset Y$ by the same lemma we used above.
Choose $V \subset \mathbf{P}^n_A$ open with $V \cap Y = j(U)$.
Then we see that (3) and (4) hold.

\medskip\noindent
Because $A$ is Noetherian we see that $X$ and $X \times_A \mathbf{P}^n_A$
are Noetherian schemes. Hence we can take $E = (b, c)_*\mathcal{O}_Z$
in this case, see Derived Categories of Schemes, Lemma
\href{https://stacks.math.columbia.edu/tag/08E8}{lemma identify pseudo coherent noetherian (Tag 08E8)}.
This finishes the proof.
\end{proof}
\end{document}
